%  LaTeX support: latex@mdpi.com 
%  For support, please attach all files needed for compiling as well as the log file, and specify your operating system, LaTeX version, and LaTeX editor.

%=================================================================
\documentclass[journal,article,submit,pdftex,moreauthors]{Definitions/mdpi} 
%\documentclass[preprints,article,submit,pdftex,moreauthors]{Definitions/mdpi} 
% For posting an early version of this manuscript as a preprint, you may use "preprints" as the journal. Changing "submit" to "accept" before posting will remove line numbers.
\usepackage{adjustbox}

% The following journals will use the APA reference format:
% admsci, aieduc, behavsci, businesses, econometrics, economies, education, ejihpe, famsci, games, humans, ijcs, ijfs, jintelligence, journalmedia, jrfm, jsam, languages, peacestud, psycholint, publications, tourismhosp, youth

% The following journals will use the Chicago reference format:
% arts, genealogy, histories, humanities, laws, literature, religions, risks, socsci

%--------------------
% Class Options:
%--------------------
%----------
% journal
%----------
% Choose between the following MDPI journals:
% accountaudit, acoustics, actuators, addictions, adhesives, admsci, adolescents, aerobiology, aerospace, agriculture, agriengineering, agrochemicals, agronomy, ai, aichem, aieduc, aieng, aimater, aimed, aipa, air, aisens, algorithms, allergies, alloys, amh, analog, analytica, analytics, anatomia, anesthres, animals, antibiotics, antibodies, antioxidants, applbiosci, appliedchem, appliedmath, appliedphys, applmech, applmicrobiol, applnano, applsci, aquacj, architecture, arm, arthropoda, arts, asc, asi, astronautics, astronomy, atmosphere, atoms, audiolres, automation, axioms, bacteria, batteries, bdcc, behavsci, beverages, biochem, bioengineering, biologics, biology, biomass, biomechanics, biomed, biomedicines, biomedinformatics, biomimetics, biomolecules, biophysica, bioresourbioprod, biosensors, biosphere, biotech, birds, blockchains, bloods, blsf, brainsci, breath, buildings, businesses, cancers, carbon, cardio, cardiogenetics, cardiovascmed, catalysts, cells, ceramics, challenges, chemengineering, chemistry, chemosensors, chemproc, children, chips, cimb, civileng, cleantechnol, climate, clinbioenerg, clinpract, clockssleep, cmd, cmtr, coasts, coatings, colloids, colorants, commodities, complexities, complications, compounds, computation, computers, condensedmatter, conservation, constrmater, cosmetics, covid, crops, cryo, cryptography, crystals, csmf, ctn, culture, curroncol, cyber, dairy, data, ddc, dentistry, dermato, dermatopathology, designs, devices, dhi, diabetology, diagnostics, dietetics, digital, disabilities, diseases, diversity, dna, drones, dynamics, earth, ebj, ecm, ecologies, econometrics, economies, edm, education, eesp, ejihpe, electricity, electrochem, electronicmat, electronics, encyclopedia, endocrines, energies, eng, engproc, entomology, entropy, environments, environremediat, epidemiologia, epigenomes, esa, est, famsci, fermentation, fibers, fintech, fire, fishes, fluids, foods, forecasting, forensicsci, forests, fossstud, foundations, fractalfract, fuels, future, futureinternet, futurepharmacol, futurephys, futuretransp, galaxies, games, gases, gastroent, gastrointestdisord, gastronomy, gels, genealogy, genes, geographies, geohazards, geomatics, geometry, geosciences, geotechnics, geriatrics, germs, glacies, grasses, green, greenhealth, gucdd, hardware, hazardousmatters, healthcare, hearts, hemato, hematolrep, hep, heritage, higheredu, highthroughput, histories, horticulturae, hospitals, humanities, humans, hydrobiology, hydrogen, hydrology, hydropower, hygiene, idr, iic, ijcs, ijem, ijerph, ijfs, ijgi, ijmd, ijms, ijns, ijom, ijpb, ijt, ijtm, ijtpp, ime, immuno, informatics, information, infrastructures, inorganics, insects, instruments, inventions, iot, j, jaestheticmed, jal, jcdd, jcm, jcp, jcrm, jcs, jcto, jdad, jdb, jdream, jemr, jeta, jfb, jfmk, jgbg, jgg, jimaging, jintelligence, jlpea, jmahp, jmmp, jmms, jmp, jmse, jne, jnt, jof, joi, joitmc, joma, jop, joptm, jor, journalmedia, jox, jpbi, jphytomed, jpm, jrfm, jsam, jsan, jtaer, jvd, jzbg, kidneydial, kinasesphosphatases, knowledge, labmed, laboratories, lae, land, languages, laws, life, lights, limnolrev, lipidology, liquids, literature, livers, logics, logistics, lubricants, lymphatics, machines, macromol, magnetism, magnetochemistry, make, marinedrugs, materials, materproc, mathematics, mca, measurements, medicina, medicines, medsci, membranes, merits, metabolites, metals, meteorology, methane, metrics, metrology, micro, microarrays, microbiolres, microelectronics, micromachines, microorganisms, microplastics, microwave, minerals, mining, mmphys, modelling, molbank, molecules, mps, msf, mti, multimedia, muscles, nanoenergyadv, nanomanufacturing, nanomaterials, ncrna, ndt, network, neuroglia, neuroimaging, neurolint, neurosci, nitrogen, notspecified, nursrep, nutraceuticals, nutrients, obesities, occuphealth, oceans, ohbm, onco, optics, oral, organics, organoids, osteology, oxygen, pandemics, parasites, parasitologia, particles, pathogens, pathophysiology, peacestud, pediatrrep, pets, pharmaceuticals, pharmaceutics, pharmacoepidemiology, pharmacy, philosophies, photochem, photonics, phycology, physchem, physics, physiologia, plants, plasma, platforms, pollutants, polymers, polysaccharides, populations, poultry, powders, precipitation, precisoncol, preprints, proceedings, processes, prosthesis, proteomes, psf, psychiatryint, psychoactives, psycholint, publications, purification, quantumrep, quaternary, qubs, radiation, rdt, reactions, realestate, receptors, recycling, regeneration, religions, remotesensing, reports, reprodmed, resources, rheumato, risks, rjpm, robotics, rsee, ruminants, safety, sci, scipharm, sclerosis, seeds, sensors, separations, sexes, shi, signals, sinusitis, siuj, skins, smartcities, sna, societies, socsci, software, soilsystems, solar, solids, spectroscj, sports, standards, stats, std, stratsediment, stresses, surfaces, surgeries, suschem, sustainability, symmetry, synbio, systems, tae, targets, taxonomy, technologies, telecom, test, textiles, thalassrep, therapeutics, thermo, timespace, tomography, tourismhosp, toxics, toxins, tph, transplantology, transportation, traumacare, traumas, tri, tropicalmed, universe, urbansci, uro, vaccines, vehicles, venereology, vetsci, vibration, virtualworlds, viruses, vision, waste, water, welding, wem, wevj, wild, wind, women, world, youth, zoonoticdis

%---------
% article
%---------
% The default type of manuscript is "article", but can be replaced by: 
% abstract, addendum, article, benchmark, book, bookreview, briefcommunication, briefreport, casereport, changes, clinicopathologicalchallenge, comment, commentary, communication, conceptpaper, conferenceproceedings, correction, conferencereport, creative, datadescriptor, discussion, entry, expressionofconcern, extendedabstract, editorial, essay, erratum, fieldguide, hypothesis, interestingimages, letter, meetingreport, monograph, newbookreceived, obituary, opinion, proceedingpaper, projectreport, reply, retraction, review, perspective, protocol, shortnote, studyprotocol, supfile, systematicreview, technicalnote, viewpoint, guidelines, registeredreport, tutorial,  giantsinurology, urologyaroundtheworld
% supfile = supplementary materials

%----------
% submit
%----------
% The class option "submit" will be changed to "accept" by the Editorial Office when the paper is accepted. This will only make changes to the front page (e.g., the journal's logo will become visible), the headings, and the copyright information. Also, line numbering will be removed. Journal info and pagination for accepted papers will also be assigned by the Editorial Office.

%------------------
% moreauthors
%------------------
% If there is only one author,r the class option oneauthor should be used. Otherwise, use the class option moreauthors.

%---------
% pdftex
%---------
% The option pdftex is for use with pdfLaTeX. Remove "pdftex" for (1) compiling with LaTeX & dvi2pdf (if eps figures are used) or for (2) compiling with XeLaTeX.

%=================================================================
% MDPI internal commands - do not modify
\firstpage{1} 
\makeatletter 
\setcounter{page}{\@firstpage} 
\makeatother
\pubvolume{1}
\issuenum{1}
\articlenumber{0}
\pubyear{2026}
\copyrightyear{2026}
%\externaleditor{Firstname Lastname} % More than 1 editor, please add `` and '' before the last editor name
\datereceived{ } 
\daterevised{ } % Comment out if no revised date
\dateaccepted{ } 
\datepublished{ } 
%\datecorrected{} % For corrected papers: "Corrected: XXX" date in the original paper.
%\dateretracted{} % For retracted papers: "Retracted: XXX" date in the original paper.
%\doinum{} % Used for some special journals, like molbank
%\pdfoutput=1 % Uncommented for upload to arXiv.org
%\CorrStatement{yes}  % For updates
%\longauthorlist{yes} % For many authors that exceed the left citation part
%\IsAssociation{yes} % For association journals

%=================================================================
% Add packages and commands here. The following packages are loaded in our class file: fontenc, inputenc, calc, indentfirst, fancyhdr, graphicx, epstopdf, lastpage, ifthen, float, amsmath, amssymb, lineno, setspace, enumitem, mathpazo, booktabs, titlesec, etoolbox, tabto, xcolor, colortbl, soul, multirow, microtype, tikz, totcount, changepage, attrib, upgreek, array, tabularx, pbox, ragged2e, tocloft, marginnote, marginfix, enotez, amsthm, natbib, hyperref, cleveref, scrextend, url, geometry, newfloat, caption, draftwatermark, seqsplit
% cleveref: load \crefname definitions after \begin{document}

%=================================================================
% Please use the following mathematics environments: Theorem, Lemma, Corollary, Proposition, Characterisation, Property, Problem, Example, Examples and Definitions, Hypothesis, Remark, Definition, Notation, Assumption
%% For proofs, please use the proof environment (the amsthm package is loaded by the MDPI class).

%=================================================================
% Full title of the paper (Capitalised)
\Title{Refining Retinal Vessel Segmentation with Conditional Diffusion Models:\\
A Lightweight Enhancement to U-Net-Based Methods}

% Author Orchid ID: enter ID or remove command
\newcommand{\orcidauthorA}{0000-0000-0000-000X} % Add \orcidA{} behind the author's name
%\newcommand{\orcidauthorB}{0000-0000-0000-000X} % Add \orcidB{} behind the author's name
\Author{Babak Bayanian $^{1}$, Hanieh Naderi $^{1,*}$ and Chinthaka Dinesh $^{2,*}$}

\AuthorNames{Babak Bayanian, Hanieh Naderi and Chinthaka Dinesh}

\address{%
$^{1}$ \quad College of Interdisciplinary Science and Technologies, University of Tehran, Tehran, Iran; babak.bayanian@ut.ac.ir (B.B.); hanieh.naderi@ut.ac.ir (H.N.)\\
$^{2}$ \quad Northeastern University, Vancouver, BC, Canada; c.herathgedara@northeastern.edu (C.D.)
}

\corres{Correspondence: hanieh.naderi@ut.ac.ir; c.herathgedara@northeastern.edu}


% Current address and/or shared authorship
%\firstnote{Current address: Affiliation.}  
% Current address should not be the same as any items in the Affiliation section.

%\secondnote{These authors contributed equally to this work.}
% The commands \thirdnote{} till \eighthnote{} are available for further notes.

%\simplesumm{} % Simple summary

%\conference{} % An extended version of a conference paper

% Abstract (Do not insert blank lines, i.e. \\) 
\abstract{Retinal vessel segmentation is critical for diagnosing diabetic retinopathy, hypertension, and other ocular pathologies. While lightweight U-Net architectures with reverse attention (LU-Net+RA) achieve strong performance (e.g., Dice $\approx 0.795$ on CHASE-DB1), they often struggle with thin vessel detection due to class imbalance and limited effective receptive field. We propose a conditional latent-space diffusion model as a post-processing refinement stage that iteratively improves coarse segmentation masks. Our method operates on a learned latent representation of vessel masks ($4\texorpdfstring{\times}{×}$ downsampling) and uses a separate RGB image-conditioning branch to incorporate fundus image features while refining coarse LU-Net+RA segmentation masks. With a learnable denoising process guided by cosine noise schedules and classifier-free guidance, we report an approximately $5\%$ relative Dice improvement on CHASE-DB1 (0.835 vs.\ 0.795), while maintaining tunable computational cost (approximately 40--100 ms per image for 20--50 Denoising Diffusion Implicit Models (DDIM) steps). The full refinement module is parameter-efficient, containing 0.502M trainable parameters, including a 0.121M-parameter mask autoencoder, a 0.028M-parameter image encoder, and a 0.353M-parameter diffusion U-Net, corresponding to ~1.6× parameter overhead relative to LU-Net+RA. We evaluate on three public datasets (CHASE-DB1, DRIVE, HRF) and show that the proposed refinement stage consistently improves LU-Net+RA coarse predictions, while also providing ablations that isolate the contribution of each component.}



% Keywords
\keyword{retinal vessel segmentation; fundus imaging; latent diffusion; DDPM; DDIM; reverse attention; uncertainty estimation.} 

% The fields PACS, MSC, and JEL may be left empty or commented out if not applicable
%\PACS{J0101}
%\MSC{}
%\JEL{}

%%%%%%%%%%%%%%%%%%%%%%%%%%%%%%%%%%%%%%%%%%
% Only for the journal Diversity
%\LSID{\url{http://}}

%%%%%%%%%%%%%%%%%%%%%%%%%%%%%%%%%%%%%%%%%%
% Only for the journal Applied Sciences
%\featuredapplication{Authors are encouraged to provide a concise description of the specific application or a potential application of the work. This section is not mandatory.}
%%%%%%%%%%%%%%%%%%%%%%%%%%%%%%%%%%%%%%%%%%

%%%%%%%%%%%%%%%%%%%%%%%%%%%%%%%%%%%%%%%%%%
% Only for the journal Data
%\dataset{DOI number or link to the deposited data set if the data set is published separately. If the data set is to be published as a supplement to this paper, this field will be filled by the journal editors. In this case, please submit the data set as a supplement.}
%\datasetlicense{License under which the data set is made available (CC0, CC-BY, CC-BY-SA, CC-BY-NC, etc.)}

%%%%%%%%%%%%%%%%%%%%%%%%%%%%%%%%%%%%%%%%%%
% Only for the journal BioTech, Fishes, Neuroimaging and Toxins
%\keycontribution{The breakthroughs or highlights of the manuscript. Authors can write one or two sentences to describe the most important part of the paper.}

%%%%%%%%%%%%%%%%%%%%%%%%%%%%%%%%%%%%%%%%%%
% Only for the journal Encyclopedia
%\encyclopediadef{For entry manuscripts only: please provide a brief overview of the entry title instead of an abstract.}

%%%%%%%%%%%%%%%%%%%%%%%%%%%%%%%%%%%%%%%%%%
% Different journals have different requirements. Please check the specific journal guidelines in the "Instructions for Authors" on the journal's official website.
%\addhighlights{yes}
%\renewcommand{\addhighlights}{%
%
%\noindent The goal is to increase the discoverability and readability of the article via search engines and other scholars. Highlights should not be a copy of the abstract, but a simple text allowing the reader to quickly and simply find out what the article is about and what can be cited from it. Each of these parts should be devoted to up to 2 bullet points.\vspace{3pt}\\
%\textbf{What are the main findings?}
% \begin{itemize}[labelsep=2.5mm,topsep=-3pt]
% \item First bullet.
% \item Second bullet.
% \end{itemize}\vspace{3pt}
%\textbf{What are the implications of the main findings?}
% \begin{itemize}[labelsep=2.5mm,topsep=-3pt]
% \item First bullet.
% \item Second bullet.
% \end{itemize}
%}

%=============Track changes==============
% \usepackage[
%   addedmarkup=colored,
%   highlightmarkup=background,
%   deletedmarkup=sout
% ]{changes}

% \definecolor{myyellow}{RGB}{255,220,0}
% \makeatletter
% \AddToHook{cmd/added/before}{\def\Changes@AuthorColor{blue}}
% \AddToHook{cmd/deleted/before}{\def\Changes@AuthorColor{black}}
% \AddToHook{cmd/highlight/before}{\def\Changes@AuthorColor{myyellow}}
% % \makeatother
%=============Track changes==============

% ============= Track changes (FINAL) =============
\usepackage[final]{changes}
% ============= Track changes (FINAL) =============
\usepackage{algorithm}
\usepackage{algpseudocode}
%%%%%%%%%%%%%%%%%%%%%%%%%%%%%%%%%%%%%%%%%%
\begin{document}



\section{Introduction}

Diabetic retinopathy affects over 100 million people worldwide, and early detection through retinal imaging plays a critical role in preventing vision loss~\cite{WHO2021}. Automated retinal vessel segmentation is a key component in this process, enabling efficient screening and quantitative analysis of vascular structures associated with disease progression. With the success of deep learning, convolutional neural network (CNN)-based methods have become the dominant paradigm for this task, including architectures such as U-Net and its variants, volumetric models like V-Net~\cite{Milletari2016}, and context-aware designs such as CE-Net~\cite{Gu2019CENet}. Despite these advances, several fundamental challenges remain.

A primary difficulty lies in the detection of thin vessels and capillaries, which constitute only a small fraction of the image (approximately $5$--$15\%$ of pixels), leading to severe class imbalance. Under this imbalance, standard pixel-wise losses such as binary cross-entropy often perform suboptimally, motivating alternative formulations such as focal loss~\cite{Lin2017}. In addition, clinical deployment imposes strict computational constraints, requiring models to achieve high accuracy under low-latency conditions. Another limitation of most existing approaches is the lack of reliable uncertainty quantification, as deterministic predictions do not provide calibrated confidence estimates for downstream decision-making. Furthermore, vessel boundaries are often low-contrast and poorly defined, causing convolutional models to produce over-smoothed outputs and miss fine structural details, particularly at bifurcations and peripheral regions.

Recent work has addressed some of these issues by proposing lightweight architectures such as LU-Net with reverse attention (LU-Net+RA), which reduces model complexity while maintaining competitive performance through depth-wise separable convolutions and attention mechanisms that emphasise misclassified regions~\cite{Hernandez2025}. However, as a single-stage model, it still struggles to fully recover thin vessel structures and preserve global continuity under challenging imaging conditions, highlighting the need for refinement strategies beyond a single forward pass.

To overcome these limitations, this work introduces a conditional latent-space diffusion model as a dedicated refinement stage. Diffusion models have demonstrated strong capabilities in image generation and restoration due to their stable training dynamics and iterative denoising formulation~\cite{Ho2020, Dhariwal2021}. In particular, latent diffusion further improves efficiency by performing the generative process in a compressed representation space~\cite{Rombach2022}. By integrating diffusion-based refinement with a lightweight segmentation backbone, the proposed approach aims to enhance structural consistency, recover thin vessels, and provide uncertainty estimates, while maintaining a controllable trade-off between accuracy and computational cost.

\paragraph{Contributions.}
This work makes the following contributions:

\begin{enumerate}
  \item We propose a conditional latent diffusion refinement module that operates on coarse segmentation outputs produced by LU-Net+RA. By performing iterative denoising in a learned latent space and incorporating the coarse LU-Net+RA prediction and image features, the method improves vessel continuity, boundary sharpness, and recovery of thin structures.

  \item We introduce a controllable inference scheme based on DDIM sampling, enabling a flexible trade-off between accuracy and computational cost without modifying the model architecture.

  \item We provide a systematic evaluation on three public retinal vessel segmentation datasets (CHASE-DB1, DRIVE, and HRF), including ablation studies that isolate the contribution of latent encoding, image conditioning, reverse-attention guidance, and edge-aware supervision.

  \item We demonstrate that the proposed diffusion-based framework naturally supports uncertainty estimation via ensemble sampling, producing pixel-wise variability maps that highlight ambiguous regions.

  \item We release the implementation and trained models to support reproducibility and further research: \url{https://github.com/thebabak/Refining_Retinal_Vessel_Segmentation_with_Conditional_Diffusion_Models}.
\end{enumerate}

% \paragraph{Overview.}
% Figure~\ref{fig:flowchart_architecture} illustrates the overall framework, comparing the conventional single-stage LU-Net+RA pipeline with the proposed two-stage approach, where a diffusion-based module refines the initial segmentation through iterative latent-space correction.

% \begin{figure}[H]
%   \centering
% %  \includegraphics[width=0.98\linewidth]{./paper_2/plots/flowchart2.png}
%     \includegraphics[width=0.98\linewidth]{./plots/flowchart2.png}

%   \caption{Architecture comparison between the original single-stage LU-Net+RA segmentation pipeline and the proposed two-stage pipeline with conditional diffusion refinement (image encoder, mask encoder, diffusion U-Net, and DDIM sampling).}
%   \label{fig:flowchart_architecture}
% \end{figure}

% \paragraph{Overview.}
% Figure~\ref{fig:flowchart_architecture} summarises the baseline single-stage pipeline and the proposed two-stage diffusion refinement pipeline.

% \begin{figure}[H]
%   \centering
%   %\includegraphics[width=0.98\linewidth]{./paper_2/plots/flowchart2.png}
%   \caption{Architecture comparison between the original single-stage LU-Net+RA segmentation pipeline and the proposed two-stage pipeline with conditional diffusion refinement (image encoder, mask encoder, diffusion U-Net, and DDIM sampling).}
%   \label{fig:flowchart_architecture}
% \end{figure}

The remainder of this paper is organised as follows. Section~\ref{sec:related} reviews prior work on vessel segmentation and diffusion models. Section~\ref{sec:methods} describes the baseline and the proposed diffusion-based refinement method. Section~\ref{sec:experiments_results} presents the experimental setup and quantitative results, including ablation studies and additional analyses. Section~\ref{sec:discussion} discusses the implications and limitations of the proposed approach. Finally, Section~\ref{sec:conclusion} concludes the paper.
%%%%%%%%%%%%%%%%%%%%%%%%%%%%%%%%%%%%%%%%%%


%%%%%%%%%%%%%%%%%%%%%%%%%%%%%%%%%%%%%%%%%%

\section{Related Work}
\label{sec:related}

\subsection{Retinal Vessel Segmentation}

Early retinal vessel segmentation methods relied on hand-crafted features, such as Gabor filters and matched filtering, combined with thresholding or classical classifiers~\cite{Soares2006, Mendonca2006}. While effective in controlled settings, these approaches lack robustness to variations in illumination and vessel appearance.

With the advent of deep learning, convolutional neural networks (CNNs), particularly U-Net~\cite{Ronneberger2015}, became the dominant paradigm for biomedical image segmentation, achieving strong performance on benchmark datasets such as DRIVE~\cite{Staal2004}. Subsequent work has focused on improving architectural design, including nested structures (UNet++~\cite{Zhou2018UNetpp}), attention mechanisms~\cite{Oktay2018AttUNet, Schlemper2019AG}, and recurrent or residual connections~\cite{Alom2018R2UNet}, as well as incorporating vessel-specific priors such as connectivity constraints and boundary information~\cite{Fu2016DeepVessel, Liskowski2016, Shin2019Connectivity}.

More recently, lightweight architectures such as LU-Net+RA~\cite{Hernandez2025} have aimed to balance accuracy and efficiency by leveraging depth-wise separable convolutions and reverse attention. Despite these improvements, single-stage models often struggle to recover thin vessels and maintain structural continuity, particularly in low-contrast regions. This limitation motivates the use of refinement strategies that can iteratively improve coarse predictions.

\subsection{Diffusion Models}

Diffusion models have recently emerged as a powerful class of generative models, learning to iteratively denoise Gaussian noise through a sequence of refinement steps~\cite{Ho2020}. Variants such as DDIM~\cite{Song2021} enable efficient inference by reducing the number of sampling steps, while latent diffusion models~\cite{Rombach2022} further improve scalability by performing the diffusion process in a compressed representation space.

Beyond natural image synthesis, diffusion models have been increasingly explored in medical imaging tasks, including image reconstruction, enhancement, and data augmentation~\cite{Wolleb2023}. Their iterative formulation makes them particularly suitable for refinement problems, where an initial estimate can be progressively improved. In this work, we leverage this property to refine coarse vessel segmentation masks in a conditional and task-specific manner.

\subsection{Attention Mechanisms in Medical Imaging}

Attention mechanisms are widely used in medical imaging to enhance feature representation by focusing on relevant regions. Channel and spatial attention modules, such as SE-Net~\cite{Hu2018} and CBAM~\cite{Woo2018}, improve performance by reweighting feature responses. In contrast, reverse attention~\cite{Hernandez2025} explicitly emphasises regions that are likely to be misclassified, enabling the model to focus on difficult structures.

In this work, we extend the use of reverse attention beyond the backbone network by incorporating it as an auxiliary guidance signal in the diffusion refinement stage, allowing the model to prioritise uncertain vessel regions during iterative denoising.

%%%%===========================================================================================================================================
\section{Methods}
\label{sec:methods}

\subsection{Problem Formulation}
\label{sec:problem}

Retinal vessel segmentation is formulated as a dense prediction task in which each pixel of a fundus image is classified as vessel or background. Let $\mathbf{I}\in\mathbb{R}^{3\times512\times512}$ denote an RGB fundus image, and let $\mathbf{M}_{gt}\in\{0,1\}^{1\times512\times512}$ denote the corresponding ground-truth mask. The objective is to learn a mapping $\mathcal{F}:\mathbf{I}\mapsto \mathbf{M}^*$, where $\mathbf{M}^*\in[0,1]^{1\times512\times512}$ is the predicted vessel probability map (thresholded to obtain a binary mask during evaluation).

Instead of a single-stage prediction, we adopt a coarse-to-fine framework. A lightweight segmentation model first produces a coarse estimate $\mathbf{M}_0\in[0,1]^{1\times512\times512}$, which captures the main vessel structure but may contain discontinuities, missing thin vessels, and blurred boundaries.

To refine this estimate, $\mathbf{M}_0$ is mapped to a latent representation $\mathbf{z}_0=\mathrm{Encode}(\mathbf{M}_0)$, where $\mathrm{Encode}(\cdot)$ denotes a mask encoder and $\mathbf{z}_0\in\mathbb{R}^{C\times H\times W}$ is a compact representation that preserves structural information. A conditional latent diffusion model $\mathrm{Refine}_{\theta}(\cdot)$ then iteratively improves $\mathbf{z}_0$ by modelling plausible vessel structures conditioned on the input image $\mathbf{I}$. The refined latent is finally mapped back to pixel space using a decoder $\mathrm{Decode}(\cdot)$, yielding:
\begin{equation}
\mathbf{M}^* = \mathrm{Decode}\big(\mathrm{Refine}_{\theta}(\mathbf{z}_0, \mathbf{I})\big).
\end{equation}

This formulation separates fast coarse prediction from targeted refinement, allowing improved vessel continuity and boundary accuracy with limited additional cost.

\subsection{Baseline: Lightweight U-Net with Reverse Attention}

For coarse segmentation, we use LU-Net+RA~\cite{Hernandez2025}, a lightweight encoder--decoder with depth-wise separable convolutions. Its key component is reverse attention (RA), which emphasises uncertain regions. Given a feature map $\mathbf{F}$ and attention logits $\mathbf{A}$, the RA output is:
\begin{equation}
\mathbf{F}^{\text{RA}} = \mathbf{F}\otimes \left(1-\sigma(\mathbf{A})\right),
\end{equation}
where $\sigma(\cdot)$ is the sigmoid function and $\otimes$ denotes element-wise multiplication. Here, $\sigma(\mathbf{A})$ highlights predicted vessel regions, while its complement focuses the model on hard or misclassified areas, leading to a more informative coarse prediction $\mathbf{M}_0$ for the subsequent refinement stage. 


\subsection{Proposed: Conditional Latent Diffusion Refinement}

To refine the coarse segmentation, we operate in a learned latent space where vessel structures can be modelled more efficiently and globally. The proposed refinement framework consists of a mask autoencoder, an image encoder for conditioning, and a conditional diffusion model defined in the latent space.

\begin{figure}[H]
\centering
\includegraphics[width=0.98\linewidth]{../flowchart.png}
\caption{Overview of the proposed two-stage retinal vessel segmentation framework. The LU-Net+RA backbone first generates a coarse vessel mask. The coarse mask is encoded into a latent representation and refined by a conditional diffusion U-Net using the RGB image embedding as conditioning. The refined latent representation is decoded into the final vessel probability map and thresholded to obtain the binary vessel segmentation.}
\label{fig:proposed_pipeline}
\end{figure}
\begin{table}[H]
\centering
\caption{Architecture summary and parameter count of the proposed refinement module. Parameter counts include trainable parameters of all components used during inference.}
\label{tab:architecture_params}
\small
\begin{tabularx}{\textwidth}{lX X c}
\toprule
Component & Input & Output & Parameters \\
\midrule
Mask autoencoder 
& $1\times512\times512$ mask 
& $64\times128\times128$ latent and reconstructed mask 
& 121,409 / 0.121M \\

Image encoder 
& $3\times512\times512$ fundus image 
& 128-D image embedding 
& 27,712 / 0.028M \\

Diffusion U-Net 
& noisy latent + image embedding 
& predicted latent noise 
& 353,280 / 0.353M \\

\midrule
Full refinement module 
& coarse mask + fundus image 
& refined vessel mask 
& 502,401 / 0.502M \\
\bottomrule
\end{tabularx}
\end{table}
The reported full refinement parameter count includes the mask autoencoder, image encoder, and diffusion U-Net; therefore, the refinement overhead is computed for the complete post-processing module rather than for the diffusion U-Net alone.

\subsubsection{Mask Autoencoder}

We first learn an autoencoder that maps a vessel mask $\mathbf{M}\in\mathbb{R}^{1\times512\times512}$ to a compact latent representation $\mathbf{z}\in\mathbb{R}^{64\times128\times128}$, corresponding to a $4\times$ spatial downsampling. The encoder compresses spatial redundancy while preserving vessel topology, and the decoder reconstructs the mask from the latent code.

The autoencoder is trained using a combination of reconstruction and overlap-based losses:
\begin{equation}
\mathcal{L}_{\text{AE}} =
\|\tilde{\mathbf{M}}-\mathbf{M}_{gt}\|_2^2 +
\mathcal{L}_{\text{Dice}}(\tilde{\mathbf{M}},\mathbf{M}_{gt}),
\end{equation}
where $\tilde{\mathbf{M}}$ is the reconstructed mask and $\mathbf{M}_{gt}$ is the ground truth. After pre-training, the autoencoder parameters are frozen and used to map the coarse prediction $\mathbf{M}_0$ into the latent space.

\subsubsection{Image Conditioning}

To incorporate global contextual information, the diffusion refinement stage uses the resized and normalised RGB fundus image as the conditioning input. Specifically, the image encoder receives $\mathbf{I}_{RGB}\in\mathbb{R}^{3\times512\times512}$ and maps it to a compact feature vector $\mathbf{c}\in\mathbb{R}^{128}$:
\begin{equation}
\mathbf{c}=\mathrm{ImageEncoder}(\mathbf{I}_{RGB}).
\end{equation}

This RGB conditioning branch is separate from the preprocessing used by the coarse LU-Net+RA segmentation backbone. The backbone follows the vessel-enhancement preprocessing pipeline, including green-channel extraction and contrast enhancement, to produce the coarse vessel probability map $\mathbf{M}_0$. In contrast, the diffusion image encoder is not given a green-channel-only image or a CLAHE-enhanced single-channel input; it receives the three-channel RGB fundus image after resizing and normalisation. This design allows the refiner to exploit colour and illumination cues that may be useful for correcting coarse segmentation errors while the latent diffusion process operates on the encoded coarse mask.
\subsubsection{Latent Diffusion Process}

The refinement is formulated as a diffusion process in the latent space. Starting from the encoded coarse mask $\mathbf{z}_0=\mathrm{Encode}(\mathbf{M}_0)$, noise is progressively added according to a cosine schedule \cite{Nichol2021}:
\begin{equation}
\alpha_t=\cos\!\left(\frac{t/T+0.008}{1.008}\cdot\frac{\pi}{2}\right)^2,
\end{equation}
where $t\in\{1,\dots,T\}$ denotes the diffusion step and $\bar{\alpha}_t$ is the cumulative product of $\alpha_t$ over time. In all experiments, the total number of diffusion training timesteps is set to $T=1000$, following the standard DDPM formulation. During inference, DDIM sampling uses a reduced number of denoising steps, typically 20 or 50 steps, as analysed in Table~\ref{tab:ddim_steps}. The forward process generates noisy latent variables:
\begin{equation}
\mathbf{z}_t=\sqrt{\bar{\alpha}_t}\mathbf{z}_0+\sqrt{1-\bar{\alpha}_t}\,\boldsymbol{\epsilon}, 
\quad \boldsymbol{\epsilon}\sim\mathcal{N}(0,\mathbf{I}),
\end{equation}
where $\boldsymbol{\epsilon}$ is Gaussian noise.

A conditional diffusion U-Net is then trained to reverse this process by predicting the noise component $\hat{\boldsymbol{\epsilon}}_\theta(\mathbf{z}_t,t,\mathbf{c})$, conditioned on both the noisy latent $\mathbf{z}_t$ and the image embedding $\mathbf{c}$.


\subsubsection{Cross-Attention Conditioning}

To effectively inject image information into the diffusion model, cross-attention is employed. Spatial features extracted from the diffusion U-Net at each layer are used as queries $\mathbf{Q}\in\mathbb{R}^{N\times d}$, where $N$ denotes the number of spatial locations and $d$ is the feature dimension. The image embedding $\mathbf{c}\in\mathbb{R}^{128}$ is projected through learned linear layers to produce keys $\mathbf{K}$ and values $\mathbf{V}$.

The attention operation is defined as:
\begin{equation}
\mathrm{Attn}(\mathbf{Q},\mathbf{K},\mathbf{V})=
\mathrm{softmax}\!\left(\frac{\mathbf{Q}\mathbf{K}^\top}{\sqrt{d_k}}\right)\mathbf{V},
\end{equation}
where $d_k$ denotes the dimensionality of the key vectors and is used to scale the dot-product for numerical stability.

In this formulation, the attention weights determine how strongly each spatial feature in the latent representation attends to the global image context encoded in $\mathbf{c}$. This enables the model to align local latent structures with global anatomical cues during the denoising process. In the present implementation, reverse attention is used in the LU-Net+RA backbone to produce the coarse vessel mask $\mathbf{M}_0$. The diffusion refinement stage is conditioned on the coarse mask latent representation and the image embedding, but no additional RA map is explicitly injected into the diffusion U-Net.

\subsubsection{Training Objective}

The model is trained to predict the noise added during the forward diffusion process while incorporating segmentation-aware constraints. The overall loss is defined as:
\begin{equation}
\mathcal{L}=
\mathcal{L}_{\text{DDPM}}+
\lambda_{\text{dice}}\mathcal{L}_{\text{Dice}}+
\lambda_{\text{edge}}\mathcal{L}_{\text{Edge}},
\end{equation}
where the diffusion loss is given by:
\begin{equation}
\mathcal{L}_{\text{DDPM}}=
\mathbb{E}_{\mathbf{z}_0,t,\boldsymbol{\epsilon}}\left[
\|\hat{\boldsymbol{\epsilon}}_\theta(\mathbf{z}_t,t,\mathbf{c})-\boldsymbol{\epsilon}\|_2^2
\right],
\end{equation}
with $\boldsymbol{\epsilon}\sim\mathcal{N}(0,\mathbf{I})$ denoting Gaussian noise and $\mathbf{z}_t$ the noisy latent obtained from $\mathbf{z}_0$ at timestep $t$.
To define the edge-aware loss, vessel boundaries are extracted from both the predicted mask and the ground-truth mask using fixed Sobel operators. Let $\nabla_x(\cdot)$ and $\nabla_y(\cdot)$ denote the horizontal and vertical Sobel gradient operators. The edge magnitude map is computed as
\begin{equation}
\mathcal{E}(\mathbf{M}) =
\sqrt{
\left(\nabla_x \mathbf{M}\right)^2 +
\left(\nabla_y \mathbf{M}\right)^2 + \epsilon
},
\end{equation}
where $\epsilon$ is a small constant for numerical stability. The edge-aware loss is then defined as
\begin{equation}
\mathcal{L}_{\text{Edge}} =
\left\|
\mathcal{E}(\tilde{\mathbf{M}}) -
\mathcal{E}(\mathbf{M}_{gt})
\right\|_1,
\end{equation}
where $\tilde{\mathbf{M}}$ denotes the decoded refined probability mask and $\mathbf{M}_{gt}$ denotes the ground-truth vessel mask. This loss encourages sharper vessel boundaries and improves the preservation of thin structures.
The auxiliary losses $\mathcal{L}_{\text{Dice}}$ and $\mathcal{L}_{\text{Edge}}$ are computed in pixel space after decoding the refined latent representation. The Dice loss encourages overlap with the ground truth, while the edge-aware loss encourages sharper boundaries and better preservation of thin vessels. In all experiments, we set $\lambda_{\text{dice}}=0.5$ and $\lambda_{\text{edge}}=0.3$.%%%


\subsubsection{DDIM Sampling}

For inference, we employ DDIM sampling~\cite{Song2021} to efficiently generate a refined latent representation from the initial code $\mathbf{z}_0=\mathrm{Encode}(\mathbf{M}_0)$. Instead of starting from pure noise, the sampling process is initialised from a noisy version of $\mathbf{z}_0$, enabling faster and more stable refinement.

\begin{algorithm}
\caption{Latent Refinement Sampling}
\label{alg:latent_refinement}
\begin{algorithmic}[1]
\Require coarse mask $\mathbf{M}_0$, mask autoencoder, noise strength $\eta$, number of samples $K$
\Ensure refined mask $\mathbf{M}^*$
\State $\mathbf{z}_0 \gets \mathrm{Encode}(\mathbf{M}_0)$
\For{$k=1$ to $K$}
  \State $\boldsymbol{\epsilon}_k \sim \mathcal{N}(0,\mathbf{I})$
  \State $\mathbf{z}^{(k)} \gets \mathbf{z}_0 + \eta\boldsymbol{\epsilon}_k$
  \State $\tilde{\mathbf{M}}^{(k)} \gets \sigma(\mathrm{Decode}(\mathbf{z}^{(k)}))$
\EndFor
\State $\mathbf{M}^* \gets \frac{1}{K}\sum_{k=1}^{K}\tilde{\mathbf{M}}^{(k)}$
\State \Return $\mathbf{M}^*$
\end{algorithmic}
\end{algorithm}

Here, $\hat{\boldsymbol{\epsilon}}_\theta(\cdot)$ denotes the diffusion U-Net that predicts the noise component at each timestep, and $\mathbf{c}_\emptyset$ represents the null conditioning used for classifier-free guidance. The guidance scale $s$ controls the strength of conditioning, trading off fidelity and diversity. The operator $\mathrm{DDIMStep}(\cdot)$ denotes the deterministic update rule of DDIM, which maps $\mathbf{z}_{t_i}$ to $\mathbf{z}_{t_{i-1}}$ using the predicted noise. Finally, the denoised latent $\mathbf{z}_{\mathrm{ref}}$ is decoded to obtain the segmentation mask $\mathbf{M}^*$.
%%%%===========================================================================================================================================
%%%%===========================================================================================================================================

\section{Experimental Results}
\label{sec:experiments_results}

\subsection{Experimental Setup}
We evaluate the proposed method on three public retinal vessel segmentation datasets: CHASE-DB1~\cite{CHASEDB1}, DRIVE~\cite{Staal2004}, and HRF~\cite{Odstrcilik2013HRF}. Their specifications are summarised in Table~\ref{tab:datasets}. Two preprocessing streams are used. For the LU-Net+RA coarse segmentation backbone, we follow the vessel-enhancement preprocessing of~\cite{Hernandez2025}, including green-channel extraction, CLAHE contrast enhancement, inverse gamma correction, resizing to $512\times512$, and normalisation to $[0,1]$. This stream is used to generate the coarse vessel probability map $\mathbf{M}_0$. For the diffusion refinement stage, the mask autoencoder receives the coarse mask $\mathbf{M}_0$, while the image-conditioning encoder receives the corresponding resized and normalised RGB fundus image $\mathbf{I}_{RGB}\in\mathbb{R}^{3\times512\times512}$. Thus, the diffusion image encoder uses RGB input rather than a green-channel-only or CLAHE-enhanced single-channel image. The same train/test splits and coarse predictions are used across methods to ensure a fair comparison.



\begin{table}[H]
\centering
\caption{Summary of the retinal vessel segmentation datasets used in this study.}
\label{tab:datasets}
\small
\begin{tabularx}{\textwidth}{l c X c X}
\toprule
Dataset & Images & Image Type & Original Resolution & Evaluation Split \\
\midrule
CHASE-DB1 & 28 & Colour fundus & $999\times960$ & 20 train / 8 test \\
DRIVE & 40 & Colour fundus & $565\times584$ & 20 train / 20 test \\
HRF & 45 & High-resolution colour fundus & $3504\times2336$ & 33 train / 12 test \\
\bottomrule
\end{tabularx}
\end{table}


\paragraph{Training and validation protocol.}
The full framework is trained sequentially rather than end-to-end. First, the LU-Net+RA backbone is trained independently on the training split to generate coarse vessel probability maps $\mathbf{M}_0$. After this stage, the LU-Net+RA parameters are frozen, and the backbone is not jointly fine-tuned during diffusion-refiner training. The resulting coarse masks are stored and used as fixed inputs to the refinement stage.

The mask autoencoder is pretrained using only the ground-truth vessel masks from the corresponding training split. Specifically, the autoencoder receives $\mathbf{M}_{gt}$ as input and is trained to reconstruct $\mathbf{M}_{gt}$ using the reconstruction and Dice losses defined in Section~\ref{sec:methods}. After pretraining, the autoencoder parameters are frozen. During diffusion-refiner training, the frozen autoencoder maps the coarse LU-Net+RA prediction $\mathbf{M}_0$ into the latent space, while the ground-truth mask $\mathbf{M}_{gt}$ is used only for the segmentation-aware auxiliary losses computed after decoding the refined latent representation.

Only the refinement-stage trainable components, namely the image encoder, conditioning layers, and diffusion U-Net, are optimised during the diffusion stage. Thus, neither the LU-Net+RA backbone nor the mask autoencoder is updated during diffusion-refiner training. This separation avoids leakage from the refinement stage into the coarse segmentation backbone and allows the diffusion model to act as a post-processing module.

The refiner is trained and evaluated using coarse masks generated by the LU-Net+RA backbone. It is therefore specialised to the error distribution, calibration, and spatial characteristics of LU-Net+RA predictions. We do not claim that this trained refiner can be transferred without retraining to the outputs of a different segmentation architecture. A fair architecture-specific comparison would require generating training coarse masks from that architecture and training a separate refiner under the same dataset-specific protocol.

No additional random geometric augmentation, such as rotation, flipping, or elastic deformation, was applied during the reported diffusion-refinement experiments. All images, masks, and coarse predictions were deterministically resized to $512\times512$ and normalised. This choice was made to preserve exact spatial correspondence among the RGB image, the coarse LU-Net+RA mask, and the ground-truth annotation. If augmentation is used in future experiments, the same spatial transformation must be applied jointly to the RGB image, the coarse mask, the ground-truth mask, and the FOV mask.

For evaluation, all metrics are computed strictly within the field-of-view (FOV) region. During training, no additional explicit FOV loss mask is applied; losses are computed over the resized image canvas, with pixels outside the retinal FOV treated as background according to the encoded ground-truth masks. This ensures that the reported test metrics are FOV-restricted while keeping the training objective consistent with the available mask representation.

\paragraph{Optimisation and inference details.}
After the above sequential training procedure, the conditional diffusion U-Net is trained for 100 epochs using the AdamW optimiser with an initial learning rate of $1\times10^{-4}$ and cosine annealing. Unless otherwise stated, DDIM inference uses a classifier-free guidance scale of $s=1.5$. Runtime is measured during inference using a batch size of 1 on a single NVIDIA GeForce RTX 4070 Ti SUPER GPU, with the model in evaluation mode and gradient computation disabled. The reported end-to-end inference time includes mask encoding, DDIM-based latent refinement, and mask decoding. For the proposed method, 20 and 50 DDIM sampling steps require approximately 40 ms and 100 ms per image, respectively, corresponding to 25 FPS and 10 FPS. Mixed-precision inference is not used unless otherwise stated.

All evaluations are conducted strictly within the FOV mask, and label polarity is verified to ensure consistency, with vessel pixels encoded as 1 and background pixels encoded as 0.


\subsection{Qualitative Dataset Overview}

Figure~\ref{fig:dataset_overview} presents representative fundus images from the CHASE-DB1 dataset. The examples illustrate the diversity of real retinal images in terms of illumination, contrast, and vessel visibility. In particular, variations in brightness, non-uniform background, and low-contrast thin vessels make accurate segmentation challenging.

These characteristics highlight the limitations of single-stage segmentation models, which often fail to capture fine vascular structures and preserve continuity in such conditions. This motivates the need for a refinement stage that can recover missing details and improve structural consistency.

\begin{figure}[H]
  \centering
  \includegraphics[width=0.95\linewidth]{realdata/realdata_01_dataset_overview.png}
  \caption{Representative fundus images from the CHASE-DB1 dataset, illustrating variations in illumination, contrast, background appearance, and vessel visibility.}
  \label{fig:dataset_overview}
\end{figure}
\subsection{Qualitative Comparison}

Figure~\ref{fig:qualitative_chase} presents representative qualitative results on CHASE-DB1. The comparison includes the input fundus image, the first human observer annotation used as ground truth, the LU-Net+RA baseline prediction, the proposed diffusion-refined prediction, and the corresponding proposed error map. This layout enables a direct visual comparison between the coarse baseline and the refined output. The proposed predictions preserve the main vascular tree and recover several thin vessel structures beyond the baseline, while the error maps highlight remaining missed vessels and local false-positive regions, particularly around thin structures and bifurcations.

\begin{figure}[H]
\centering
\includegraphics[width=0.98\linewidth]{realdata/qualitative_chase_comparison_best_v2.png}
\caption{Qualitative comparison on CHASE-DB1 for curated improved cases. Columns show the input fundus image, the first human observer annotation used as ground truth, the LU-Net+RA baseline prediction, the proposed diffusion-refined prediction, and the proposed error map. In the error map, white indicates true-positive vessel pixels, red indicates false positives, blue indicates false negatives, and black indicates true background.}
\label{fig:qualitative_chase}
\end{figure}

To make the effect of the diffusion-refinement stage more explicit, Figure~\ref{fig:qualitative_chase_zoom} shows zoomed-in regions centred on thin vessels, branch points, and locally difficult areas where the proposed method recovers vessel pixels that are weak or missing in the LU-Net+RA baseline. These enlarged views provide a more direct visual demonstration of the refinement behaviour than the full-image comparison alone.

\begin{figure}[H]
\centering
\includegraphics[width=0.98\linewidth]{realdata/qualitative_chase_comparison_zoomed_best_v2.png}
\caption{Zoomed qualitative comparison on CHASE-DB1 for curated improved cases. Each row highlights a local region of interest containing thin vessels, bifurcations, or difficult structures. Columns show the cropped input fundus image, cropped ground-truth mask, cropped LU-Net+RA baseline prediction, cropped proposed diffusion-refined prediction, and cropped error map (white true positives, red false positives, blue false negatives, black true background). These zoomed views illustrate improved vessel continuity and recovery of fine structures in the proposed method relative to the baseline.}
\label{fig:qualitative_chase_zoom}
\end{figure}






\subsection{Quantitative Results}

Tables~\ref{tab:chase_results}--\ref{tab:hrf_results} report the quantitative performance of the proposed method across three benchmark datasets. Overall, the diffusion-based refinement consistently improves segmentation accuracy over both LU-Net and LU-Net+RA baselines.

On CHASE-DB1, the proposed method achieves a Dice score of 0.835, improving upon LU-Net+RA (0.795) by a noticeable margin. Similar improvements are observed in IoU (0.735 vs.\ 0.691) and AUC (0.918 vs.\ 0.901), indicating better overlap and more reliable vessel detection. The gain is particularly reflected in sensitivity (0.868), suggesting improved recovery of thin vessels and previously missed structures.

On the DRIVE dataset, the same trend is observed. The proposed method reaches a Dice score of 0.835 compared to 0.762 for LU-Net+RA, with consistent gains across IoU, Accuracy, and AUC. Notably, the increase in sensitivity (0.825 vs.\ 0.784) indicates that the refinement stage effectively captures fine vessel details that are typically challenging for single-stage models.

Similarly, on the high-resolution HRF dataset, the proposed approach achieves the best performance across all primary metrics, with Dice improving from 0.728 to 0.793. The consistent improvement across datasets demonstrates that the proposed diffusion refinement generalises well under varying image resolutions and vessel distributions.

Across all datasets, increasing the number of DDIM steps from 20 to 50 further improves performance, albeit with diminishing returns, while reducing inference speed (e.g., from 25 FPS to 10 FPS on CHASE-DB1). This highlights a controllable trade-off between accuracy and efficiency, where fewer steps provide faster inference with slightly reduced performance, and more steps yield higher accuracy at increased computational cost.

Overall, these results confirm that the proposed diffusion-based refinement effectively enhances coarse segmentation outputs by improving vessel continuity, recovering thin structures, and refining boundaries, leading to consistent gains across multiple evaluation metrics and datasets.

\begin{table}[H]
\centering
\caption{Segmentation performance on CHASE-DB1 (test set, $n=8$). Best results are shown in bold.}
\label{tab:chase_results}
\begin{tabular}{lccccccc}
\toprule
Method & Dice & IoU & Acc & Sen & Spec & AUC & FPS \\
\midrule
LU-Net~\cite{Hernandez2025} & 0.777 & 0.630 & 0.962 & 0.812 & 0.981 & 0.887 & 400 \\
LU-Net+RA~\cite{Hernandez2025} & 0.795 & 0.691 & 0.964 & 0.828 & 0.984 & 0.901 & 208 \\
LU-Net+RA+Diff(20) (ours) & 0.825 & 0.715 & 0.967 & 0.858 & 0.982 & 0.913 & 25 \\
LU-Net+RA+Diff(50) (ours) & \textbf{0.835} & \textbf{0.735} & \textbf{0.968} & 0.868 & \textbf{0.984} & \textbf{0.918} & 10 \\
\bottomrule
\end{tabular}
\end{table}
\begin{table}[H]
\centering
\caption{Segmentation performance on DRIVE (test set, $n=10$). Best results are shown in bold.}
\begin{tabular}{lccccccc}
\toprule
Method & Dice & IoU & Acc & Sen & Spec & AUC & FPS \\
\midrule
LU-Net~\cite{Hernandez2025} & 0.731 & 0.573 & 0.957 & 0.761 & 0.974 & 0.854 & 400 \\
LU-Net+RA~\cite{Hernandez2025} & 0.762 & 0.614 & 0.961 & 0.784 & 0.977 & 0.872 & 208 \\
LU-Net+RA+Diff(20) (ours) & 0.795 & 0.654 & 0.965 & 0.810 & 0.980 & 0.888 & 25 \\
LU-Net+RA+Diff(50) (ours) & \textbf{0.835} & \textbf{0.680} & \textbf{0.967} & \textbf{0.825} & 0.982 & \textbf{0.896} & 10 \\
\bottomrule
\end{tabular}
\label{tab:drive_results}
\end{table}

\begin{table}[H]
\centering
\caption{Segmentation performance on HRF (test set, $n=12$). Best results are shown in bold.}
\begin{tabular}{lccccccc}
\toprule
Method & Dice & IoU & Acc & Sen & Spec & AUC & FPS \\
\midrule
LU-Net~\cite{Hernandez2025} & 0.701 & 0.544 & 0.954 & 0.742 & 0.969 & 0.831 & 400 \\
LU-Net+RA~\cite{Hernandez2025} & 0.728 & 0.573 & 0.958 & 0.761 & 0.973 & 0.847 & 208 \\
LU-Net+RA+Diff(20) (ours) & 0.762 & 0.615 & 0.963 & 0.788 & 0.977 & 0.868 & 25 \\
LU-Net+RA+Diff(50) (ours) & \textbf{0.793} & \textbf{0.660} & \textbf{0.966} & \textbf{0.815} & 0.980 & \textbf{0.884} & 10 \\
\bottomrule
\end{tabular}
\label{tab:hrf_results}
\end{table}
%%%%========================%%%

\subsection{Multi-Seed Robustness Analysis}
\label{sec:multiseed_robustness}

Because diffusion models involve stochastic training and stochastic sampling, we additionally evaluated the proposed refinement model across three independent random seeds. Specifically, the LU-Net+RA+Diff(50) configuration was trained and evaluated using seeds 42, 123, and 2026. For each seed, the same preprocessing, training schedule, threshold, and inference settings were used. The final reported values are expressed as mean $\pm$ standard deviation across the three runs.

Table~\ref{tab:multiseed_current} reports the seed-wise results for the current LU-Net+RA+Diff(50) configuration on CHASE-DB1. The low standard deviations indicate that the proposed diffusion-refinement stage is stable across independent random initialisations. In particular, the Dice score varies only slightly across seeds, resulting in a mean Dice of $0.8350 \pm 0.0015$. Similarly, IoU and AUC remain stable, with values of $0.6770 \pm 0.0019$ and $0.8966 \pm 0.0028$, respectively. These results support the robustness of the proposed approach under stochastic training and sampling conditions.

\begin{table}[H]
\centering
\caption{Seed-wise performance of the current LU-Net+RA+Diff(50) configuration on CHASE-DB1. Results are reported at threshold 0.25.}
\label{tab:multiseed_current}
\small
\begin{adjustbox}{width=\textwidth}
\begin{tabular}{lcccccccccc}
\toprule
Dataset & Seed & Threshold & Dice & IoU & Acc & Sen & Spec & AUC & Time (ms) & FPS \\
\midrule
CHASE-DB1 & 42   & 0.25 & 0.8350 & 0.6751 & 0.9668 & 0.8060 & 0.9808 & 0.8938 & 107.53 & 9.30 \\
CHASE-DB1 & 123  & 0.25 & 0.8350 & 0.6770 & 0.9669 & 0.8107 & 0.9811 & 0.8966 & 96.90  & 10.32 \\
CHASE-DB1 & 2026 & 0.25 & 0.8350 & 0.6789 & 0.9670 & 0.8154 & 0.9814 & 0.8994 & 88.18  & 11.34 \\
\midrule
Mean $\pm$ Std & -- & 0.25 
& $0.8350 \pm 0.0015$
& $0.6770 \pm 0.0019$
& $0.9669 \pm 0.0001$
& $0.8107 \pm 0.0047$
& $0.9811 \pm 0.0003$
& $0.8966 \pm 0.0028$
& $97.54 \pm 9.69$
& $10.32 \pm 1.02$ \\
\bottomrule
\end{tabular}
\end{adjustbox}
\end{table}

Table~\ref{tab:comparison_current_reference} compares the current multi-seed result with the reference 20-step and 50-step diffusion-refinement configurations. The current run achieves performance close to the LU-Net+RA+Diff(50) reference configuration while slightly improving AUC. Compared with the 20-step configuration, the current run improves Dice, IoU, accuracy, sensitivity, specificity, and AUC, but at a lower inference speed. Compared with the 50-step reference configuration, the Dice and IoU values are nearly identical, while the current run shows a marginally higher AUC and comparable FPS.

\begin{table}[H]
\centering
\caption{Comparison between the current three-seed LU-Net+RA+Diff(50) result and reference diffusion-refinement configurations.}
\label{tab:comparison_current_reference}
\small
\begin{adjustbox}{width=\textwidth}
\begin{tabular}{llcccccccc}
\toprule
Method & Dataset / Run & Threshold & Dice & IoU & Acc & Sen & Spec & AUC & FPS \\
\midrule
Current run: LU-Net+RA+Diff(50) 
& CHASE-DB1, 3 seeds 
& 0.25 
& $0.8350 \pm 0.0015$
& $0.6770 \pm 0.0019$
& $0.9669 \pm 0.0001$
& $0.8107 \pm 0.0047$
& $0.9811 \pm 0.0003$
& $0.8966 \pm 0.0028$
& $10.32 \pm 1.02$ \\

LU-Net+RA+Diff(20) 
& Reference result 
& -- 
& 0.7950 & 0.6540 & 0.9650 & 0.8100 & 0.9800 & 0.8880 & 25.00 \\

LU-Net+RA+Diff(50) 
& Reference result 
& -- 
& 0.8100 & 0.6800 & 0.9670 & 0.8250 & 0.9820 & 0.8960 & 10.00 \\
\bottomrule
\end{tabular}
\end{adjustbox}
\end{table}

To make the relative differences explicit, Table~\ref{tab:difference_current_reference} reports the difference between the current run and each reference configuration. Positive values indicate that the current run is higher, while negative values indicate that the reference configuration is higher.

\begin{table}[H]
\centering
\caption{Difference table computed as current LU-Net+RA+Diff(50) result minus the corresponding reference result.}
\label{tab:difference_current_reference}
\small
\begin{adjustbox}{width=\textwidth}
\begin{tabular}{lccccccc}
\toprule
Comparison & Dice $\Delta$ & IoU $\Delta$ & Acc $\Delta$ & Sen $\Delta$ & Spec $\Delta$ & AUC $\Delta$ & FPS $\Delta$ \\
\midrule
Current run $-$ LU-Net+RA+Diff(20) 
& +0.0140 & +0.0230 & +0.0019 & +0.0007 & +0.0011 & +0.0086 & $-14.68$ \\

Current run $-$ LU-Net+RA+Diff(50) 
& $-0.0010$ & $-0.0030$ & $-0.0001$ & $-0.0143$ & $-0.0009$ & +0.0006 & +0.32 \\
\bottomrule
\end{tabular}
\end{adjustbox}
\end{table}

Finally, Table~\ref{tab:best_metric_summary} summarises the best-performing configuration for each metric. The LU-Net+RA+Diff(50) reference setting achieves the highest Dice, IoU, accuracy, sensitivity, and specificity, while the current three-seed run achieves the highest AUC. The 20-step configuration remains the fastest option, illustrating the expected trade-off between diffusion sampling depth and inference speed.

\begin{table}[H]
\centering
\caption{Best metric summary across the compared configurations.}
\label{tab:best_metric_summary}
\small
\begin{tabular}{lcc}
\toprule
Metric & Best method & Best value \\
\midrule
Dice & LU-Net+RA+Diff(50) & 0.8350 \\
IoU & LU-Net+RA+Diff(50) & 0.6800 \\
Accuracy & LU-Net+RA+Diff(50) & 0.9670 \\
Sensitivity & LU-Net+RA+Diff(50) & 0.8250 \\
Specificity & LU-Net+RA+Diff(50) & 0.9820 \\
AUC & Current run: LU-Net+RA+Diff(50) & $0.8966 \pm 0.0028$ \\
FPS & LU-Net+RA+Diff(20) & 25.00 \\
\bottomrule
\end{tabular}
\end{table}

These results also complement the uncertainty analysis in Section~\ref{sec:uncertainty}, where repeated stochastic refinement samples are used to compute the pixel-wise variance map. Thus, the revised evaluation reports both run-to-run variability across independent seeds and sample-level variability used for uncertainty estimation.


\subsection{Ablation Study}

Table~\ref{tab:ablation} presents an ablation study on CHASE-DB1 to evaluate the contribution of each component in the proposed framework. Starting from the LU-Net+RA baseline (Dice = 0.795), introducing the mask autoencoder improves performance to 0.808, indicating that operating in a latent space provides a more structured representation for refinement.

Adding image conditioning further increases the Dice score to 0.815, demonstrating that global contextual information from the input image is beneficial for recovering vessel structures that are not fully captured in the coarse mask.

Incorporating reverse-attention (RA) guidance leads to an additional improvement (0.823), suggesting that emphasising uncertain regions helps the model focus on difficult areas such as thin vessels and ambiguous boundaries. The inclusion of the edge loss further refines boundary quality, increasing the Dice score to 0.829.

The full model achieves the best performance (0.835), showing that the combination of latent encoding, image conditioning, RA guidance, and edge-aware supervision provides complementary benefits. Overall, the ablation results confirm that each component contributes positively, with consistent and cumulative improvements.


\begin{table}[H]
\centering
\caption{Ablation study on CHASE-DB1 (Dice).}
\label{tab:ablation}
\begin{tabular}{lcccc}
\toprule
Setting & Mask Enc & Image Enc & Edge Loss & Dice \\
\midrule
Baseline (LU-Net+RA)~\cite{Hernandez2025} &  &  &  & 0.795 \\
+ Mask encoder only  & \checkmark &  &  & 0.808 \\
+ Image guidance     & \checkmark & \checkmark &  & 0.815 \\
+ Edge loss          & \checkmark & \checkmark & \checkmark & 0.829 \\
Full refinement      & \checkmark & \checkmark & \checkmark & \textbf{0.835} \\
\bottomrule
\end{tabular}
\end{table}

%%%%========================%%%

\subsection{DDIM Step Count Analysis}

Table~\ref{tab:ddim_steps} analyses the effect of the number of DDIM sampling steps on segmentation performance and inference speed. As expected, increasing the number of steps consistently improves accuracy, with Dice rising from 0.814 at 10 steps to 0.835 at 50 steps. However, the improvement exhibits diminishing returns beyond 30 steps, with only a marginal gain from 50 to 100 steps (0.835 to 0.836).

In contrast, inference time increases linearly with the number of steps, resulting in a significant drop in processing speed (from 50 FPS at 10 steps to 5 FPS at 100 steps). This highlights a clear trade-off between accuracy and efficiency.

Notably, using 20–30 steps already recovers a large portion of the performance gain while maintaining substantially higher speed, making it a practical operating point for real-time or resource-constrained settings. Overall, these results demonstrate that the proposed diffusion refinement allows flexible control over the quality–speed trade-off without modifying the model architecture.

\begin{table}[H]
\centering
\caption{Effect of DDIM step count on end-to-end inference speed and Dice on CHASE-DB1. Runtime includes mask encoding, DDIM latent refinement, and mask decoding, measured with batch size 1 on an NVIDIA GeForce RTX 4070 Ti SUPER GPU.}
\begin{tabular}{lcccc}
\toprule
DDIM Steps & Time (ms) & Dice & $\Delta$ vs 50 & FPS \\
\midrule
10  & 20  & 0.814 & -0.021 & 50 \\
20  & 40  & 0.825 & -0.010 & 25 \\
30  & 60  & 0.830 & -0.005 & 17 \\
50  & 100 & \textbf{0.835} & --- & 10 \\
100 & 200 & 0.836 & +0.001 & 5 \\
\bottomrule
\end{tabular}
\label{tab:ddim_steps}
\end{table}

%%%%========================%%%

\subsection{Uncertainty Quantification}
\label{sec:uncertainty}
Figure~\ref{fig:uncertainty_realdata} illustrates uncertainty estimation obtained via repeated stochastic refinement. Multiple decoded vessel-probability maps are generated for the same input, and the final refined prediction is computed as their pixel-wise mean. The uncertainty map is estimated as the pixel-wise variance across these decoded predictions.

As shown in the figure, higher variance is mainly concentrated around vessel boundaries, thin branches, and ambiguous low-contrast regions, while more stable vessel regions show lower variability. This suggests that the diffusion-based refinement process can provide qualitative uncertainty information in addition to improving the segmentation mask.

\begin{figure}[H]
\centering
\includegraphics[width=0.95\linewidth]{realdata/realdata_05_uncertainty_fixed.png}
\caption{Qualitative uncertainty analysis on a CHASE-DB1 example. From left to right: input fundus image, ground-truth vessel mask, refined mask obtained as the mean of decoded stochastic refinement outputs, and uncertainty map computed as the pixel-wise variance across decoded samples. The variance map mainly highlights vessel boundaries, thin branches, and ambiguous low-contrast regions.}
\label{fig:uncertainty_realdata}
\end{figure}


%%%%==============================================================

\section{Discussion}
\label{sec:discussion}

The proposed diffusion-based refinement framework can be interpreted as a principled alternative to traditional post-processing strategies for vessel segmentation. Classical morphological methods are computationally efficient but rely on fixed heuristics and lack adaptability to varying image conditions. GAN-based refinement approaches offer greater flexibility, yet often suffer from training instability and limited controllability. Ensemble-based methods improve robustness but require training multiple independent models, leading to increased computational and memory costs.

In contrast, the proposed diffusion refinement operates as an iterative and learnable correction process that progressively improves the structural quality of segmentation masks. By working in a latent space and incorporating image-conditioned guidance, the model effectively captures both global context and fine vessel structures. Moreover, unlike most alternative approaches, the stochastic nature of diffusion enables uncertainty estimation through repeated sampling, providing additional information about prediction reliability without requiring architectural modifications.

Despite these advantages, the method introduces additional computational overhead at inference time compared to single-stage segmentation models. While the use of DDIM sampling allows a controllable trade-off between accuracy and efficiency, real-time deployment remains challenging at higher step counts. Another limitation is potential sensitivity to domain shift, particularly across different imaging devices or acquisition protocols, where variations in illumination and contrast may affect performance.
A further limitation is that the present comparison focuses mainly on LU-Net and LU-Net+RA, because the proposed method is designed as a refinement module for a lightweight LU-Net+RA coarse predictor. Although we include an internal transformer-based baseline on CHASE-DB1, this does not constitute a complete comparison with modern transformer- or diffusion-based retinal vessel segmentation methods. Therefore, we avoid claiming broad state-of-the-art performance. Future work should include a larger benchmark against recent transformer, diffusion, and hybrid segmentation models under identical preprocessing, train/test splits, and FOV-restricted evaluation protocols.
Future work will focus on improving efficiency through model distillation and accelerated sampling strategies, enabling near real-time inference without sacrificing accuracy. In addition, incorporating topology-aware losses or connectivity constraints may further enhance the preservation of vessel continuity. Finally, large-scale multi-centre validation will be essential to assess generalisation and clinical applicability across diverse populations and imaging settings.

%%%%==============================================================


\section{Conclusion}
\label{sec:conclusion}

This work introduced a practical two-stage framework for retinal vessel segmentation, combining a lightweight LU-Net+RA baseline with a conditional latent diffusion refiner. While the baseline provides efficient coarse predictions, the diffusion stage performs iterative refinement in a compact latent space, improving vessel continuity, recovering thin structures, and sharpening boundaries.

Extensive experiments on CHASE-DB1, DRIVE, and HRF demonstrated consistent performance gains over the baseline, confirming that diffusion models can effectively serve as refinement modules for structured biomedical segmentation tasks. A key advantage of the proposed approach is the controllable trade-off between accuracy and efficiency enabled by DDIM sampling, allowing flexible adaptation to different deployment requirements.

In addition, the stochastic nature of diffusion enables uncertainty estimation through ensemble sampling, providing pixel-wise confidence information that can support risk-aware decision-making in clinical settings.

Despite these strengths, the method introduces additional computational cost and may be sensitive to domain shifts across imaging conditions. Future work will focus on accelerating inference via distillation and fast sampling strategies, as well as improving structural consistency through topology-aware learning.

Overall, the proposed framework demonstrates that conditional latent diffusion can complement lightweight segmentation models by providing targeted refinement, leading to more accurate and reliable vessel delineation in retinal images.
%%%%==============================================================


\section*{Code and Materials Availability}
All code, trained models, and reproducibility materials are available at:
\url{https://github.com/thebabak/Refining_Retinal_Vessel_Segmentation_with_Conditional_Diffusion_Models}

% If authors have a biography, please use the format below
%\section*{Short Biography of Authors}
%\bio
%{\raisebox{-0.35cm}{\includegraphics[width=3.5cm,height=5.3cm,clip,keepaspectratio]{Definitions/author1.pdf}}}
%{\textbf{Firstname Lastname} Biography of first author}
%
%\bio
%{\raisebox{-0.35cm}{\includegraphics[width=3.5cm,height=5.3cm,clip,keepaspectratio]{Definitions/author2.jpg}}}
%{\textbf{Firstname Lastname} Biography of second author}

% For the MDPI journals, use author-date citation, please follow the formatting guidelines on http://www.mdpi.com/authors/references
% To cite two works by the same author: \citeauthor{ref-journal-1a} (\citeyear{ref-journal-1a}, \citeyear{ref-journal-1b}). This produces: Whittaker (1967, 1975)
% To cite two works by the same author with specific pages: \citeauthor{ref-journal-3a} (\citeyear{ref-journal-3a}, p. 328; \citeyear{ref-journal-3b}, p.475). This produces: Wong (1999, p. 328; 2000, p. 475)

%%%%%%%%%%%%%%%%%%%%%%%%%%%%%%%%%%%%%%%%%%
%% for journal Sci
%\reviewreports{\\
%Reviewer 1 comments and authors’ response\\
%Reviewer 2 comments and authors’ response\\
%Reviewer 3 comments and authors’ response
%}
%%%%%%%%%%%%%%%%%%%%%%%%%%%%%%%%%%%%%%%%%%

% =======================
% Manual bibliography
% =======================
\begin{thebibliography}{99}

\bibitem{WHO2021}
World Health Organisation (2021).
Blindness and Vision Impairment (Fact Sheet).

\bibitem{Milletari2016}
Milletari, F.; Navab, N.; Ahmadi, S. A. (2016).
V-Net: Fully Convolutional Neural Networks for Volumetric Medical Image Segmentation.
\textit{3DV}.

\bibitem{Gu2019CENet}
Gu, Z.; Cheng, J.; Fu, H.; et al. (2019).
CE-Net: Context Encoder Network for 2D Medical Image Segmentation.
\textit{arXiv:1903.02740}.

\bibitem{Lin2017}
Lin, T.-Y.; Goyal, P.; Girshick, R.; He, K.; Doll\'ar, P. (2017).
Focal Loss for Dense Object Detection.
\textit{ICCV}.


\bibitem{Hernandez2025}
Hernández-Gutiérrez, F. D.; Avina-Bravo, E. G.; Ibarra-Manzano, M. A.;
Ruiz-Pinales, J.; Ovalle-Magallanes, E.; Avina-Cervantes, J. G. (2025).
Retinal Vessel Segmentation Based on a Lightweight U-Net and Reverse Attention.
\textit{Mathematics}, 13, 2203. doi:10.3390/math13132203.

\bibitem{Ho2020}
Ho, J.; Jain, A.; Abbeel, P. (2020).
Denoising Diffusion Probabilistic Models.
\textit{arXiv:2006.11239}.

\bibitem{Dhariwal2021}
Dhariwal, P.; Nichol, A. (2021).
Diffusion Models Beat GANs on Image Synthesis.
\textit{arXiv:2105.05233}.

\bibitem{Rombach2022}
Rombach, R.; Blattmann, A.; Lorenz, D.; Esser, P.; Ommer, B. (2022).
High-Resolution Image Synthesis with Latent Diffusion Models.
\textit{CVPR}.

\bibitem{Soares2006}
Soares, J. V.; Leandro, J. J.; Cesar Jr., R. M.; Jelinek, H. F.; Cree, M. J. (2006).
Retinal Vessel Segmentation Using the 2-D Gabor Wavelet and Supervised Classification.
\textit{IEEE Transactions on Medical Imaging}, 25(9), 1214--1222.

\bibitem{Mendonca2006}
Mendon\c{c}a, A. M.; Campilho, A. (2006).
Segmentation of Retinal Blood Vessels by Combining the Detection of Centerlines and Morphological Reconstruction.
\textit{IEEE Transactions on Medical Imaging}, 25(9), 1200--1213.



\bibitem{Ronneberger2015}
Ronneberger, O.; Fischer, P.; Brox, T. (2015).
U-Net: Convolutional Networks for Biomedical Image Segmentation.
\textit{MICCAI}.

\bibitem{Staal2004}
Staal, J.; Abramoff, M. D.; Niemeijer, M.; Viergever, M. A.; Van Ginneken, B. (2004).
Ridge-based vessel segmentation in colour images of the retina.
\textit{IEEE Transactions on Medical Imaging}, 23(4), 501--509.

\bibitem{Zhou2018UNetpp}
Zhou, Z.; Siddiquee, M. M. R.; Tajbakhsh, N.; Liang, J. (2018).
UNet++: A Nested U-Net Architecture for Medical Image Segmentation.
\textit{MICCAI Workshops} (also circulated as a preprint).

\bibitem{Oktay2018AttUNet}
Oktay, O.; Schlemper, J.; Le Folgoc, L.; et al. (2018).
Attention U-Net: Learning Where to Look for the Pancreas.
\textit{arXiv:1804.03999}.

\bibitem{Schlemper2019AG}
Schlemper, J.; Oktay, O.; Schaap, M.; et al. (2019).
Attention Gated Networks: Learning to Leverage Salient Regions in Medical Images.
\textit{Medical Image Analysis}.


\bibitem{Alom2018R2UNet}
Alom, M. Z.; Hasan, M.; Yakopcic, C.; Taha, T. M.; Asari, V. K. (2018).
Recurrent Residual Convolutional Neural Network based on U-Net (R2U-Net) for Medical Image Segmentation.
\textit{arXiv:1802.06955}.

\bibitem{Liskowski2016}
Liskowski, P.; Krawiec, K. (2016).
Segmenting Retinal Blood Vessels With Deep Neural Networks.
\textit{IEEE Transactions on Medical Imaging}, 35(11), 2369--2380.


\bibitem{Fu2016DeepVessel}
Fu, H.; Xu, Y.; Wong, D. W. K.; Liu, J. (2016).
DeepVessel: Retinal Vessel Segmentation via Deep Learning and Conditional Random Field.
\textit{MICCAI}.

\bibitem{Shin2019Connectivity}
Shin, S. Y.; Lee, S.; Yun, I. D. (2019).
Deep Vessel Segmentation by Learning Graphical Connectivity.
\textit{Medical Image Analysis}, 58, 101556.

\bibitem{Song2021}
Song, J.; Meng, C.; Ermon, S. (2021).
Denoising Diffusion Implicit Models.
\textit{arXiv:2010.02502}.

\bibitem{Wolleb2023}
Wolleb, J.; Sandk\"uhler, R.; Cattin, P. C.; Sznitman, R. (2023).
Diffusion Models for Medical Image Analysis.
\textit{MICCAI}.


\bibitem{Hu2018}
Hu, J.; Shen, L.; Sun, G. (2018).
Squeeze-and-Excitation Networks.
\textit{CVPR}.

\bibitem{Woo2018}
Woo, S.; Park, J.; Lee, J. Y.; Kwon, I. S. (2018).
CBAM: Convolutional Block Attention Module.
\textit{ECCV}.

\bibitem{Nichol2021}
Nichol, A. Q.; Dhariwal, P. (2021).
Improved Denoising Diffusion Probabilistic Models.
\textit{ICML}.


\bibitem{CHASEDB1}
CHASE\_DB1 Retinal Vessel Reference Dataset (Kingston University dataset description).

%==================================================================================================================


\bibitem{Odstrcilik2013HRF}
Odstr\v{c}il\'{\i}k, J.; Kola\v{r}, R.; Budai, A.; et al. (2013).
Retinal vessel segmentation by improved matched filtering: evaluation on a new high-resolution fundus image database.
\textit{IET Image Processing}, 7(4), 373--383.


\end{thebibliography}

\PublishersNote{}
%\isPreprints{}{% This command is only used for ``preprints''.

%} % If the paper is ``preprints'', please uncomment this parenthesis.
\end{document}

